\section{Preliminaries and Scope}
\label{sec:preliminaries-scope}

\subsection{Abstract rewriting}
\label{subsec:abstract-rewriting}

An \emph{abstract rewriting system} is a pair $(X,\to)$ consisting of a set $X$ of states and a binary relation ${\to} \subseteq X \times X$, the one-step rewrite relation. No finiteness is assumed. We write $\to^{*}$ for the reflexive--transitive closure of $\to$, and say that $b$ is \emph{reachable} from $a$, or is a \emph{descendant} of $a$, when $a \to^{*} b$. The following notions are standard \cite{BaaderNipkow1998,Huet1980,Terese2003}.

\begin{definition}[Basic notions]
\label{def:basic-notions}
Let $(X,\to)$ be an abstract rewriting system.
\begin{enumerate}
    \item States $b,c \in X$ are \emph{joinable}, written $b \joinable c$, if there is a $d \in X$ with $b \to^{*} d$ and $c \to^{*} d$.
    \item A state $n$ is a \emph{normal form} if there is no $m$ with $n \to m$.
    \item The system is \emph{locally confluent} if $a \to b$ and $a \to c$ imply $b \joinable c$.
    \item The system is \emph{confluent} if $a \to^{*} b$ and $a \to^{*} c$ imply $b \joinable c$.
    \item The system is \emph{normalizing} if every state reaches some normal form, and \emph{terminating} if there is no infinite rewrite sequence $a_0 \to a_1 \to a_2 \to \cdots$.
    \item A state $a$ has \emph{unique reachable normal forms} if any two normal forms reachable from $a$ are equal. This asserts that there is \emph{at most} one; existence is a separate matter.
\end{enumerate}
\end{definition}

Termination implies normalization, but not conversely; and neither confluence nor uniqueness of normal forms implies either of them. When $X$ is finite, the system is the same thing as a finite directed graph, possibly with self-loops, and termination amounts to the absence of directed cycles.

In string and term rewriting, $X$ is a set of words or first-order terms and $\to$ is generated by a set of rules applied at any position and, for terms, under any substitution. These settings have additional structure, notably critical pairs, which we recall in \Cref{sec:critical-pairs-curvature}.

\subsection{Place in the Six Birds program}
\label{subsec:six-birds-placement}

\begin{remark}[Six Birds primitives]
\label{rem:six-birds-primitives}
This paper uses the closure vocabulary introduced in \cite{Tsiokos2026SixBirdsFoundations}, whose six primitives are operator rewrite (P1), constraints or feasibility gating (P2), route mismatch or holonomy (P3), staging (P4), packaging (P5), and accounting or audit (P6). The present paper is primarily about P3: its subject is local route mismatch in rewriting systems. P5 and P6 play supporting roles, in packaging rewriting data into reduction 2-complexes and in keeping every defect explicit in machine-readable records. The mathematics below does not depend on this vocabulary.
\end{remark}

Within that program, the paper is a mathematical companion to \cite{Tsiokos2026CountStone}, which treats route mismatch as a quantitative defect, and it draws on the finite-state and coarse-graining background of \cite{Tsiokos2026LayStone}. Neither is needed to read what follows.

\subsection{Three kinds of evidence}
\label{subsec:scope-evidence}

The paper combines three kinds of evidence, and we keep them separate throughout.

\begin{enumerate}
    \item \emph{Proved results.} The statements in \Cref{sec:core-theorem} are theorems about arbitrary abstract rewriting systems, with paper proofs and machine-checked Lean proofs. \Cref{prop:bare-graph-not-diagnostic} is proved by explicit counterexamples.
    \item \emph{Cited classical results.} The critical-pair lemma of Knuth--Bendix and Huet, which connects critical pairs to local confluence in string and term rewriting, is used as stated in the literature. We neither re-prove nor mechanize it.
    \item \emph{Computations.} An exhaustive enumeration of small finite systems, a search for smallest counterexamples, and a set of curated string and term-rewriting examples. Statements resting on these computations concern only the systems examined. Where a computation would need an unbounded search, as with term rewriting from a start term, a verdict is reported only after the set of reachable terms has been explored completely.
\end{enumerate}

The paper does not claim a new confluence criterion for term rewriting, nor a holonomy theory in the sense of a group-valued connection. Its contribution is a precise local dictionary, the global theorems that follow from it, and a computational test of both.
